\documentclass[english,project]{buwthesis}
\begin{document}
\chapter{Literature Review}


\section{Literature Review} \label{sec:literature_review}

Autonomous vehicles integrate perception, navigation, and control to enable independent operation, as exemplified by early ROS-based systems like TurtleBot~\cite{quigley2009ros}. These systems rely on software engineering to combine technologies such as Simultaneous Localization and Mapping (SLAM), path planning, and motion control. While much of the literature targets advanced implementations, our project explores a basic autonomous vehicle using ROS2 Humble, a Raspberry Pi 4, and low-cost sensors (Section 3). This review examines key works relevant to our setup, highlighting their approaches and identifying areas less emphasized in formal research that our efforts aim to investigate.

% \subsection{SLAM Techniques}
SLAM enables a vehicle to simultaneously map its surroundings and estimate its position, a fundamental capability for autonomy. Thrun et al. (2005) introduced probabilistic SLAM with Extended Kalman Filters (EKF), a widely adopted method that fuses sensor data effectively but demands significant computational resources, posing challenges for embedded systems like our Raspberry Pi 4~\cite{thrun2005probabilistic}. For LiDAR-based approaches, Zhang and Singh (2014) developed LOAM, which achieves precise mapping with sensors like our RPLIDAR A1, though its high hardware requirements limit its use on low-cost platforms~\cite{zhang2014loam}. Seeking a more practical solution, Macenski (2020) created the SLAM Toolbox for ROS2, designed to reduce computational load and support lightweight LiDAR systems~\cite{macenski2020slamtoolbox}. Its compatibility with our setup makes it a promising approach, though adapting SLAM to resource-constrained hardware remains underexplored in the literature, an area our project seeks to investigate further.

% \subsection{Navigation and Path Planning}
Navigation and path planning enable a vehicle to move autonomously by determining optimal routes and avoiding obstacles, a key requirement for independent operation. Marder-Eppstein et al. (2010) established the ROS Navigation Stack, which uses A* for global path planning and the Dynamic Window Approach (DWA) for local adjustments, offering a robust framework but demanding precise control often beyond our basic hardware~\cite{marder2010navigation}. Building on this, Macenski et al. (2022) introduced Navigation2 (Nav2) for ROS2, enhancing modularity with behavior trees and retaining A* and DWA, making it suitable for our project\textquotesingle s software stack (Section 3.1.3)~\cite{macenski2022nav2}. Its design supports flexible integration with low-cost systems like our Raspberry Pi 4, though adapting it to coarse actuators, such as our DC motor steering with limited precision (Section 4.1.1), poses practical challenges. The cited research often assumes precise hardware for effective navigation, leaving challenges in applying tools like Navono2 to basic setups with coarse actuators, such as ours, less emphasized an aspect our project seeks to investigate.

% \subsection{Control Mechanisms}
Control mechanisms ensure a vehicle\textquotesingle s stable and accurate motion, a critical aspect of autonomous operation. \AA strom and H\"agglund
 (2006) established foundational principles for Proportional-Integral-Derivative (PID) controllers, a widely used method for regulating speed and steering, though tuning them effectively requires precise actuators often beyond our basic setup~\cite{astrom2006pid}. In the context of ROS2, the ros2\_control framework (2022) integrates PID control with navigation outputs like Nav2, offering a flexible solution for managing motor commands on platforms such as our Raspberry Pi 4 (Section 3.1.4)~\cite{ros2humble2022control}. Its design supports low-cost systems, yet adapting it to coarse hardware, such as our DC motor steering with limited precision (Section 4.1.1), presents tuning challenges. The cited research often assumes refined actuators for seamless control, leaving difficulties in applying tools like ros2\_control to basic setups with simpler hardware, such as ours, less highlighted an aspect our project seeks to examine.

% \subsection{Synthesis and Project Relevance}
Early systems like TurtleBot~\cite{quigley2009ros} demonstrate the integration of SLAM, navigation, and control within ROS frameworks, providing a foundation for our work. Our project leverages lightweight tools SLAM Toolbox, Nav2, and ros2\_control tailored to a basic setup with a Raspberry Pi 4 and low-cost sensors (Section 3). While the cited research focuses on robust implementations with advanced hardware, it less frequently addresses the practical challenges of adapting these technologies to simpler configurations, such as ours with coarse actuators and limited resources. Our efforts aim to explore these underexplored aspects, offering initial insights into their application on resource-constrained platforms.

\end{document}